\lecture{Schedulazione Classica}{18~-~02~-~2026}
%
%
\section{Fondamenti della Schedulazione}
La schedulazione è il processo decisionale che riguarda l'allocazione di risorse nel tempo per l'esecuzione di compiti.
Risolve problemi decisionali di allocazione ottimale di risorse scarse (il tempo) alle attività, di breve e brevissimo periodo, a livello di singola macchina e con capacità finita.

Per affrontare questi problemi intervengono:
\begin{itemize}
    \item \textbf{Job}: insieme di operazioni da svolgere su una risorsa;
    \item \textbf{Task}: singola operazione di un job svolta su una risorsa (attività);
    \item \textbf{Resource}: entità che effettua l'operazione (macchina);
\end{itemize}
e la cui soluzione può essere:
\begin{itemize}
    \item \textbf{Sequenza}: ordine con cui i job devono essere eseguiti su ciascuna macchina;
    \item \textbf{Schedule}: sequenza comprendente gli istanti di inizio e fine della lavorazione dei job;
    \item \textbf{Dispatching list}: schedule in cui è indicata la risorsa da utilizzare;
\end{itemize}
ottenuta in modo:
\begin{itemize}
    \item \textbf{Statico}: tutti i job sono disponibili all'inizio del problema;
    \item \textbf{Dinamico}: nuovi job possono arrivare nel tempo.
\end{itemize}

%
\paragraph{Ipotesi generali} La stessa risorsa non può effettuare più di un'operazione per volta, ogni job deve essere lavorato da una risorsa per volta e le precedenze devono essere rispettate.

%
\paragraph{Ipotesi Semplificative} I modelli classici si basano su assunzioni rigide che semplificano la realtà:
\begin{itemize}
    \item Tutti i job sono disponibili all'inizio del problema ($ t=0 $).
    \item I tempi di processo sono deterministici.
    \item I tempi di processo non dipendono dalla sequenza (no setups).
    \item Le risorse sono sempre disponibili (no failures).
    \item Le risorse lavorano un unico job per volta.
    \item Ogni job visita una risorsa al più una volta.
    \item Non è possibile interrompere la lavorazione di un job.
    \item Non vi è cancellazione di job.
    \item I lotti sono indivisibili.
    \item È possibile avere WIP.
\end{itemize}

%
\subsection{Risorse}

\begin{gather*}
    \text{Risorse} = \left\{ m_{1},m_{2}, \dots, m_{i}, \dots, m_{m} \right\} \\
    \text{Risorse}~i = 1, \dots, m
\end{gather*}

Parametri:
\begin{itemize}
    \item $ 1 $ Single Machine: tutti i job vengono eseguiti sull'unica risorsa;
    \item $ P_{m} $ Paralle Machine: i job possono essere eseguiti su una delle $ m $ risorse (identiche o meno) disponibili;
    \item $ F_{m} $ Flow shop: i job devono essere eseguiti nello stesso ordine sulle $ m $ macchine disponibili
          \begin{itemize}
              \item \emph{Skip flow shop}: alcuni job vengono eseguiti su una sola macchina;
              \item \emph{Re-entrant flow shop}: alcuni job presentano dei cicli;
          \end{itemize}
    \item $ O_{m} $ Open shop: i job devono essere eseguiti sulle $ m $ macchine disponibili;
    \item $ J_{m} $ Job shop: ciascun job deve essere eseguito su un diverso sottoinsieme delle $ m $ macchine disponibili.
\end{itemize}

%
\subsection{Misure}
Per valutare la bontà di una schedulazione, utilizziamo diverse metriche temporali:

\begin{multicols*}{2}
    %
    \paragraph{Completition time} Istante in cui il job $ j $ termina.
    \[
        C_{j}
    \]

    \paragraph{Lateness} Differenza tra il Completion time e la due date del job $ j $. Se positiva indica ritardo, se negativa anticipo.
    \[
        L_{j}= C_{j} - d_{j}
    \]

    %
    \paragraph{Tardiness} Indica il ritardo, coincide con la Lateness se questa è positiva, è nulla in caso contrario.
    \[
        T_{j} = \max(0,L_{j})
    \]

    %
    \paragraph{Earliness} Indica l'anticipo, coincide (in valore assoluto) con la Lateness se questa è negativa, è nulla in caso contrario.
    \[
        E_{j} = \max(0,-L_{j})
    \]

    %
    \paragraph{Flow time} Differenza tra il Completion time e il release time del job $ j $.
    \[
        F_{j} = C_{j}-r_{j}
    \]

    %
    \paragraph{Makespan} Istante in cui termina l'ultima operazione dell'ultimo job. Coincide con il Completion time dell'ultimo job.
    \[
        C_{\max}
    \]

    %
    \paragraph{Idleness} Differenza tra il Makespan e la somma dei process time dei job eseguiti sulla macchina $ i $.
    \[
        I_{i} = C_{\max}-\textstyle\sum_{j}p_{ij}
    \]

    \begin{note}
        Alcune misure sono equivalenti. Ad esempio, minimizzare la lateness massima ($ L_{\max} $) equivale a minimizzare la tardiness massima ($ T_{\max} $) se esiste almeno un job in ritardo.
    \end{note}
\end{multicols*}

%
\paragraph{Massimo tempo di completamento o Makespan}
\[
    C_{\max} = \text{max}_{j}C_{j}
\]
Misura rispetto all'istante $ 0 $ del tempo necessario a completare le attività.

%
\paragraph{Massima Lateness}
\[
    L_{\max} = \text{max}_{j}L_{j}
\]
Misura del ritardo del job che è maggiormente in ritardo rispetto alla sua data di consegna (se negativa misura l'anticipo del job completato con minore anticipo rispetto alla sua data di consegna).

%
\paragraph{Massima Tardiness}
\[
    T_{\max} = \max(0,L_{\max})
\]
Misura del ritardo del job che è maggiormente in ritardo rispetto alla sua data di consegna.

%
\paragraph{Tempo di completamento totale}
\[
    \textstyle\sum_{j}C_{j} \cdot (\textstyle\sum _{j}w_{j}C_{j})
\]


%
\section{Notazione di Graham}
I problemi di schedulazione sono formalizzati tramite la notazione:
\[
    \alpha / \beta / \gamma
\]
dove:
\begin{gather*}
    \alpha = \text{sistema e numero di macchine}; \\
    \beta = \text{eventuali caratteristiche particolari}; \\
    \gamma = \text{funzione obiettivo}
\end{gather*}

\begin{example}
    Può essere un esempio della notazione di Graham:
    \[
        1 / / \textstyle \sum_{j}C_{j}
    \]
    Problema su macchina singola che richiede la minimizzazione del tempo di completamento totale (Risolto con la regola SPT).
\end{example}

\begin{example}
    Un'altro esempio:
    \[
        F2 / / C_{\max}
    \]
    Problema di flow shop su due macchine che richiede la minimizzazione del makespan (Risolto con l'algoritmo di Johnson).
\end{example}

\begin{example}
    Infine:
    \[
        1 / \text{prec} / L_{\max}
    \]
    Problema su macchina singola con vincoli di precedenza tra i job che richiede la minimizzazione della massima lateness (Risolto con l'algoritmo di Lawler).
\end{example}


%
\section{Regole di priorità statistiche}
Queste regole, dette ``One-shot'', ordinano i job all'inizio e non cambiano durante l'esecuzione. Sono fondamentali perché forniscono soluzioni ottime per specifici problemi su macchina singola ($ 1 / / \dots $).

%
\paragraph{SPT (Shortest Process Time)} La priorità di un job è inversamente proporzionale al suo tempo di processo. Sequenzia i job in ordine non decrescente di $ p_{j} $. Fornisce la soluzione ottima del problema $ 1 / / \sum_{j} C_{j} $
%
\paragraph{WSPT (Weighted Shortest Process Time)} Analoga alla SPT tiene conto dei pesi dei job definendo la priorità come rapporto $ w_{j} / p_{j} $. Ordina i job secondo valori non crescenti di $ w_{j} / p_{j} $. Fornisce la soluzione ottima del problema $ 1 / / \sum_{j} w_{j}C_{j} $

%
\paragraph{LPT (Longest Process Time)} La priorità di un job è proporzionale al suo tempo di processo. Sequenzia i job in ordine non crescente di $ p_{j} $. Fornisce una soluzione non ottima del problema $ P_{m} / / C_{\max} $ bilanciando i carichi di lavoro. Si usa quando i buffer di una linea sono pieni consentendo alle macchine a valle di svuotarli

%
\section{Algoritmi risolutivi}
%
\subsection{Algo: Johnson}
Risolve in modo \emph{ottimo} il problema $ F2 / / C_{\max} $.
\begin{enumerate}[leftmargin=2.1cm, label=Passo \arabic*.]
    \item Porre $ k=1 $ e $ l=n $ numero di job;
    \item Cercare $ \min\{ p_{ij} \} $:
          \begin{itemize}
              \item Se $ \min_{i,j} \{ p_{ij} \} = \min_{j} \{ p_{1j} \}$ mettere $ j $ in posizione $ k $ e $ k=k+1 $;
              \item Se $ \min_{i,j} \{ p_{ij} \} = \min_{j} \{ p_{2j} \}$ mettere $ j $ in posizione $ l $ e $ l=l-1 $;
          \end{itemize}
    \item Se la sequenza non è completa rifare il passo 2.
\end{enumerate}

%
\subsection{Algo: Campbell}
Fornisce una soluzione non ottima al problema $ F_{m} / / C_{\max} $, è un metodo euristico.
\begin{enumerate}[leftmargin=2.1cm, label=Passo \arabic*.]
    \item Per $ k=1, \dots , m-1 $ creare le macchine:
          \[
              M1^{k}=\left\{\bigcup_{l=1}^{k}M_l\right\},\quad P1_{j}^{k}=\sum_{l=1}^{k}p_{jl},\quad M2^{k}=\left\{\bigcup_{l=k+1}^{m}M_l\right\},\quad P2_{j}^{k}=\sum_{l=k+1}^{m}p_{jl}.
          \]
          Applicare l'algoritmo di Johnson alle macchine $ M 1^{k} $ e $ M 2^{k} $ e calcolare $ C_{\max}^{k} $;
    \item Identificare la soluzione che fornisce il minimo makespan:
          \[
              C_{\text{sol}} = \min_{k=1, \dots, m-1} \left\{ C_{\max}^{k} \right\}
          \]
\end{enumerate}

%
\subsection{Algo: Lawler}
Fornisce una soluzione al problema $ 1 / \text{prec} / L_{\max} $, tenendo conto dei vincoli di precedenza tra i job.

\todo[inline]{Aggiungi il resto \dots}
